\documentclass[11pt,a4paper]{article}

\usepackage[utf8]{inputenc}
\usepackage[T1]{fontenc}
\usepackage{lmodern}
\usepackage[margin=1in]{geometry}
\usepackage{amsmath,amssymb,amsthm}
\usepackage{array}
\usepackage{booktabs}
\usepackage{xcolor}
\usepackage{hyperref}
\usepackage{microtype}

\hypersetup{
  colorlinks=true,
  linkcolor=blue!60!black,
  citecolor=blue!60!black,
  urlcolor=blue!60!black,
}

\theoremstyle{definition}
\newtheorem{definition}{Definition}
\newtheorem{example}[definition]{Example}

\theoremstyle{plain}
\newtheorem{lemma}[definition]{Lemma}
\newtheorem{proposition}[definition]{Proposition}
\newtheorem{theorem}[definition]{Theorem}
\newtheorem{corollary}[definition]{Corollary}

\theoremstyle{remark}
\newtheorem{remark}[definition]{Remark}

\newcommand{\Tv}{\mathbf{T}}
\newcommand{\Fv}{\mathbf{F}}
\newcommand{\Uv}{\mathbf{U}}
\newcommand{\Vals}{\{\Tv, \Fv, \Uv\}}
\newcommand{\Vh}{\widehat{V}}
\newcommand{\frm}{(l, \theta, o)}
\newcommand{\entails}{\vDash}

\title{Location, Time, and Observer Based Three-Valued Logic:\\
A Minimal Formal System Where Truth is Never Absolute}
\author{Norbert Nopper}
\date{\today}

\begin{document}
\maketitle

\begin{abstract}
\noindent
We present \textbf{LTO-K3}, a minimal formal system in which the truth value
of a proposition is never absolute but is always evaluated with respect to
three explicit parameters: a \emph{location}, a \emph{time}, and an
\emph{observer}. The system admits exactly three truth values---$\Tv$~(true),
$\Fv$~(false), and $\Uv$~(unknown)---composed via Kleene's strong three-valued
connectives. We give a formal syntax, a compositional point-wise semantics
$\Vh : \mathcal{L} \times L \times \Theta \times O \to \Vals$, two consequence
relations (local and universal), and prove a small body of results:
information-monotonicity, conservativity over classical logic, locality,
absence of $\Tv$-tautologies (a known property of $K_3$ that we re-derive in
the indexed setting), and the De~Morgan laws. We are careful to distinguish
what is genuinely new (the \emph{packaging}: location and observer as
non-optional, concrete parameters; a deliberate refusal to introduce modal
operators) from what is inherited from established work (Kleene
semantics~\cite{kleene1952}, modal/temporal/epistemic
indexing~\cite{prior1957,hintikka1962,kripke1963}, three-valued spatio-temporal
logics in verification). The contribution is therefore a careful synthesis,
not a discovery, and we name its limitations explicitly.
\end{abstract}

\noindent\textbf{Keywords:} three-valued logic, Kleene logic, indexed
semantics, temporal logic, epistemic logic, contextual semantics, partial
information.

\section{Introduction}
Classical two-valued logic forces every proposition into one of two states:
\emph{true} or \emph{false}. In many practical settings---verification of
distributed systems, knowledge representation, robotics, legal reasoning, data
quality---this dichotomy is too coarse. Truth often depends on \emph{where} an
evaluation is performed, \emph{when} it is performed, and \emph{who} performs
it, and in many cases the evaluator simply does not know.

We introduce a minimal logic, \textbf{LTO-K3} (Location, Time, Observer;
Kleene three-valued), that takes these dependencies as primitive. Every
atomic proposition is evaluated under an explicit tuple $\frm$ of location,
time, and observer, and may take one of three values: true, false, or
unknown. Compound formulas are evaluated point-wise via Kleene's strong
three-valued tables.

The system is intentionally austere. It contains no modal operators, no
accessibility relations, no transition systems, and no quantification over
locations, times, or observers. Its purpose is to serve as a compact,
unambiguous core on which richer domain-specific layers (synchronization,
trust, frame alignment, modal operators) may later be built.

\paragraph{Contributions.}
\begin{enumerate}
  \item A formal syntax and compositional point-wise semantics for a
  three-valued indexed logic over $(L, \Theta, O)$ frames (Sections~\ref{sec:syntax}--\ref{sec:semantics}).
  \item Two consequence relations---local at a frame and universal across
  frames---based on the designated value $\Tv$ (Section~\ref{sec:consequence}).
  \item A small body of meta-theoretic results (Section~\ref{sec:properties}):
  information-monotonicity, locality, conservativity over classical logic,
  De~Morgan laws, and the absence of $\Tv$-tautologies in the unindexed
  fragment.
  \item An explicit comparison with the closest neighbours: Kleene $K_3$,
  {\L}ukasiewicz $\text{\L}_3$, Belnap's $\mathbf{FOUR}$, Priest's LP,
  Bochvar $B_3$, hybrid logic, and multi-agent epistemic temporal logic
  (Section~\ref{sec:related}).
  \item A frank limitations section (Section~\ref{sec:limitations}).
\end{enumerate}

\paragraph{What is and is not new.}
None of the parts is new. Three-valued logics go back to
{\L}ukasiewicz~\cite{lukasiewicz1920} and Kleene~\cite{kleene1952}. Temporal
indexing goes back to Prior~\cite{prior1957}; possible-worlds indexing to
Kripke~\cite{kripke1963}; agent indexing to Hintikka~\cite{hintikka1962} and,
in computer science, to Fagin~et~al.~\cite{fagin1995}. What is new is the
\emph{packaging}: making location, time, and observer all simultaneously
non-optional, concrete, first-class parameters of a single valuation
function, and deliberately refusing to lift them with modal operators. We
state this honestly throughout.

\section{Preliminaries: Kleene's Strong Three-Valued Logic}
\label{sec:prelim}

Kleene's strong three-valued logic $K_3$~\cite{kleene1952} extends classical
two-valued logic with a third value $\Uv$ (``unknown'' or ``undefined'') and
defines the connectives so that $\Uv$ propagates conservatively: a connective
returns $\Tv$ (resp.\ $\Fv$) iff every way of resolving each $\Uv$ to $\Tv$
or $\Fv$ would yield $\Tv$ (resp.\ $\Fv$); otherwise it returns $\Uv$. This
``supervaluation-like'' character is what makes $K_3$ the natural choice
when $\Uv$ is read epistemically as ``not yet known''. The truth tables are:

\begin{center}
$\begin{array}{c|ccc}
\land & \Tv & \Uv & \Fv \\ \hline
\Tv & \Tv & \Uv & \Fv \\
\Uv & \Uv & \Uv & \Fv \\
\Fv & \Fv & \Fv & \Fv
\end{array}$\qquad
$\begin{array}{c|ccc}
\lor & \Tv & \Uv & \Fv \\ \hline
\Tv & \Tv & \Tv & \Tv \\
\Uv & \Tv & \Uv & \Uv \\
\Fv & \Tv & \Uv & \Fv
\end{array}$\qquad
$\begin{array}{c|c}
p & \lnot p \\ \hline
\Tv & \Fv \\
\Uv & \Uv \\
\Fv & \Tv
\end{array}$\qquad
$\begin{array}{c|ccc}
\to & \Tv & \Uv & \Fv \\ \hline
\Tv & \Tv & \Uv & \Fv \\
\Uv & \Tv & \Uv & \Uv \\
\Fv & \Tv & \Tv & \Tv
\end{array}$
\end{center}

Material implication is defined as $\varphi \to \psi \;:=\; \lnot \varphi
\lor \psi$. We use this strong-Kleene implication throughout; for the
contrast with {\L}ukasiewicz's $\text{\L}_3$ (where $\Uv \to \Uv = \Tv$) see
Section~\ref{sec:related}.

\paragraph{Information ordering.}
Define the partial order $\sqsubseteq$ on $\Vals$ by $\Uv \sqsubseteq \Tv$,
$\Uv \sqsubseteq \Fv$, and $x \sqsubseteq x$ for all $x$; $\Tv$ and $\Fv$
are incomparable. Intuitively, $x \sqsubseteq y$ means ``$y$ contains at
least as much information as $x$.'' All Kleene connectives are
\emph{monotone} in $\sqsubseteq$ component-wise---resolving an $\Uv$ to
$\Tv$ or $\Fv$ never reverses an already-determined truth value. This
monotonicity is the formal reason $K_3$ is appropriate when $\Uv$ denotes
epistemic uncertainty.

\section{Syntax}
\label{sec:syntax}

Fix a countably infinite set $\mathrm{AP}$ of \emph{atomic propositions}
$p, q, r, \ldots$. The language $\mathcal{L}$ is generated by the grammar
\[
\varphi, \psi \;::=\; p \;\mid\; \lnot\varphi \;\mid\; (\varphi \land \psi)
\;\mid\; (\varphi \lor \psi) \;\mid\; (\varphi \to \psi),
\qquad p \in \mathrm{AP}.
\]
We omit parentheses under the usual precedence
$\lnot > \land > \lor > {\to}$. The language has no constants, no
biconditional, and no quantifiers. The set of atomic propositions occurring
in $\varphi$ is written $\mathrm{AP}(\varphi)$.

\begin{remark}
The biconditional $\varphi \leftrightarrow \psi$ and the quantifiers
$\forall, \exists$ are deliberately omitted. The biconditional can be
defined as $(\varphi \to \psi) \land (\psi \to \varphi)$ but its $K_3$
behaviour is non-standard (it is not the equivalence relation classically
expected) and we prefer to omit it from the core. Quantifiers would force a
domain commitment and are orthogonal to the indexing concern of this paper.
Crucially, the language contains \emph{no modal operators} indexed over
$L$, $\Theta$, or $O$ (no $\Box_l$, no $\mathsf{F}\theta$, no
$K_o\varphi$). This is the essential design choice: indices appear in the
meta-language (the semantics), not in the object language.
\end{remark}

\section{Semantics}
\label{sec:semantics}

\subsection{Frames and models}

\begin{definition}[Frame]
A \emph{frame} is a triple $\mathcal{F} = (L, \Theta, O)$ where $L$, $\Theta$,
and $O$ are non-empty sets of \emph{locations}, \emph{times}, and
\emph{observers}, respectively. We call $\frm \in L \times \Theta \times O$
a \emph{point} of the frame.
\end{definition}

We impose no structure on $L$, $\Theta$, or $O$. In particular $\Theta$ is
not required to be ordered, $L$ not required to carry a topology, and $O$
not required to carry an information-state relation. Any such structure is
a domain-specific extension; see Section~\ref{sec:limitations}.

\begin{definition}[Model]
An \emph{LTO-K3 model} is a pair $\mathcal{M} = (\mathcal{F}, V)$ where
$\mathcal{F} = (L, \Theta, O)$ is a frame and
\[
V : \mathrm{AP} \times L \times \Theta \times O \;\longrightarrow\; \Vals
\]
is a \emph{base valuation} that assigns a truth value to each atomic
proposition at each point. Note that $V$ is total: every atomic proposition
has a value (possibly $\Uv$) at every point.
\end{definition}

\subsection{Compositional valuation}

The base valuation $V$ extends uniquely to a valuation $\Vh$ on all formulas
by structural recursion using the Kleene tables of Section~\ref{sec:prelim}.

\begin{definition}[Extended valuation]
Let $\mathcal{M} = (\mathcal{F}, V)$ be a model with $\mathcal{F} =
(L, \Theta, O)$. The \emph{extended valuation}
$\Vh : \mathcal{L} \times L \times \Theta \times O \to \Vals$ is defined by
recursion on $\varphi$:
\begin{align*}
\Vh(p, l, \theta, o) &\;=\; V(p, l, \theta, o), \quad p \in \mathrm{AP}, \\
\Vh(\lnot\varphi, l, \theta, o) &\;=\; \lnot_{K_3} \Vh(\varphi, l, \theta, o), \\
\Vh(\varphi \land \psi, l, \theta, o) &\;=\; \Vh(\varphi, l, \theta, o) \land_{K_3} \Vh(\psi, l, \theta, o), \\
\Vh(\varphi \lor \psi, l, \theta, o) &\;=\; \Vh(\varphi, l, \theta, o) \lor_{K_3} \Vh(\psi, l, \theta, o), \\
\Vh(\varphi \to \psi, l, \theta, o) &\;=\; \Vh(\varphi, l, \theta, o) \to_{K_3} \Vh(\psi, l, \theta, o),
\end{align*}
where $\lnot_{K_3}, \land_{K_3}, \lor_{K_3}, \to_{K_3}$ are the connectives
of Section~\ref{sec:prelim}.
\end{definition}

This is the entire semantics. Crucially, the recursion is \emph{point-wise}:
both sides of every binary connective are evaluated at \emph{the same}
point $\frm$. There is no syntactic mechanism in $\mathcal{L}$ for
referring to a different point.

\begin{remark}[Cross-frame composition]
A natural question is how to combine $\Vh(\varphi, l_1, \theta_1, o_1)$
with $\Vh(\psi, l_2, \theta_2, o_2)$. The answer in LTO-K3 is: \emph{the
language does not express such combinations}. There is no formula whose
two sub-formulas are evaluated at different points; the recursion above
forces a single $\frm$ throughout. Defining cross-frame operators would
require explicit bridging rules (synchronization, trust, observer
alignment) that are domain-specific. Any such extension is outside the
core system.
\end{remark}

\section{Consequence Relations}
\label{sec:consequence}

We take $\Tv$ as the sole designated value. (Designating $\{\Tv, \Uv\}$
gives Priest's LP~\cite{priest1979}, a different logic; see
Section~\ref{sec:related}.)

\begin{definition}[Local satisfaction]
Let $\mathcal{M} = (\mathcal{F}, V)$ be a model and $\frm$ a point. We say
$\varphi$ is \emph{satisfied at} $\frm$ in $\mathcal{M}$, written
$\mathcal{M}, \frm \entails \varphi$, iff $\Vh(\varphi, l, \theta, o) = \Tv$.
\end{definition}

\begin{definition}[Consequence]
Let $\Gamma \subseteq \mathcal{L}$ and $\varphi \in \mathcal{L}$.
\begin{itemize}
  \item \emph{Local entailment at} $\frm$ in $\mathcal{M}$:
  $\Gamma \entails_{\mathcal{M}, \frm} \varphi$ iff
  $\mathcal{M}, \frm \entails \gamma$ for every $\gamma \in \Gamma$ implies
  $\mathcal{M}, \frm \entails \varphi$.
  \item \emph{Universal entailment}: $\Gamma \entails \varphi$ iff
  $\Gamma \entails_{\mathcal{M}, \frm} \varphi$ for every model
  $\mathcal{M}$ and every point $\frm$ of $\mathcal{M}$.
\end{itemize}
\end{definition}

These are the standard $K_3$-style consequence relations, lifted point-wise
over the frame.

\section{Properties}
\label{sec:properties}

We collect the basic meta-theoretic properties of LTO-K3. The proofs are
short because the semantics is point-wise: at each fixed $\frm$ the system
collapses to ordinary $K_3$, so every $K_3$ identity lifts.

\begin{lemma}[Locality]\label{lem:locality}
For every model $\mathcal{M}$, every formula $\varphi$, and every point
$\frm$, the value $\Vh(\varphi, l, \theta, o)$ depends only on the
restriction of $V$ to $\mathrm{AP}(\varphi) \times \{l\} \times \{\theta\}
\times \{o\}$.
\end{lemma}

\begin{proof}
Structural induction on $\varphi$. The base case ($\varphi = p$) is
immediate by the definition of $\Vh$ on atoms. For each compound case
$\varphi = \lnot\psi$, $\psi \land \chi$, $\psi \lor \chi$, $\psi \to \chi$,
the value at $\frm$ is a function (the corresponding Kleene connective) of
$\Vh(\psi, l, \theta, o)$ and possibly $\Vh(\chi, l, \theta, o)$, which by
the induction hypothesis depend only on $V$ restricted as claimed.
$\mathrm{AP}(\varphi)$ is the union of the $\mathrm{AP}$ of its
sub-formulas, so the restriction propagates. \qedhere
\end{proof}

\begin{proposition}[Information-monotonicity]\label{prop:mono}
Let $V_1, V_2$ be base valuations on the same frame with
$V_1(p, l, \theta, o) \sqsubseteq V_2(p, l, \theta, o)$ for every atom and
every point. Then for every formula $\varphi$ and every point $\frm$,
$\Vh_1(\varphi, l, \theta, o) \sqsubseteq \Vh_2(\varphi, l, \theta, o)$.
\end{proposition}

\begin{proof}
Structural induction, using that $\lnot_{K_3}, \land_{K_3}, \lor_{K_3},
\to_{K_3}$ are each monotone in $\sqsubseteq$ in each argument (a standard
$K_3$ result, verified by inspection of the truth tables).
\end{proof}

\begin{corollary}[Conservativity over classical logic]\label{cor:cons}
Let $\mathcal{M}$ be a model in which $V$ never takes the value $\Uv$.
Then for every formula $\varphi$ and every point $\frm$,
$\Vh(\varphi, l, \theta, o) \in \{\Tv, \Fv\}$, and the resulting value
coincides with the classical (two-valued) valuation of $\varphi$ under the
same atomic assignment.
\end{corollary}

\begin{proof}
By Proposition~\ref{prop:mono} and inspection of the Kleene tables: when
neither input is $\Uv$, every Kleene connective returns the classical
value.
\end{proof}

\begin{proposition}[De Morgan laws]\label{prop:demorgan}
For all formulas $\varphi, \psi$ and all $\frm$,
\begin{align*}
\Vh(\lnot(\varphi \land \psi), l, \theta, o) &= \Vh(\lnot\varphi \lor \lnot\psi, l, \theta, o), \\
\Vh(\lnot(\varphi \lor \psi), l, \theta, o) &= \Vh(\lnot\varphi \land \lnot\psi, l, \theta, o).
\end{align*}
\end{proposition}

\begin{proof}
By case analysis on $(\Vh(\varphi), \Vh(\psi)) \in \Vals \times \Vals$
(nine cases per identity). Each case checks directly against the truth
tables of Section~\ref{sec:prelim}.
\end{proof}

\begin{proposition}[Double negation, idempotence, commutativity, associativity]\label{prop:basic}
For all $\varphi, \psi, \chi$ and all $\frm$, $\Vh$ satisfies:
$\lnot\lnot\varphi \equiv \varphi$;\;
$\varphi \land \varphi \equiv \varphi$;\;
$\varphi \lor \varphi \equiv \varphi$;\;
$\varphi \land \psi \equiv \psi \land \varphi$;\;
$\varphi \lor \psi \equiv \psi \lor \varphi$;\;
$(\varphi \land \psi) \land \chi \equiv \varphi \land (\psi \land \chi)$;\;
$(\varphi \lor \psi) \lor \chi \equiv \varphi \lor (\psi \lor \chi)$,
where $\equiv$ means equality of $\Vh$-values at every point of every
model.
\end{proposition}

\begin{proof}
Each is a finite truth-table check, lifted point-wise via
Lemma~\ref{lem:locality}.
\end{proof}

\begin{theorem}[No universal $\Tv$-tautologies]\label{thm:notaut}
There is no formula $\varphi \in \mathcal{L}$ such that
$\entails \varphi$ (i.e., $\Vh(\varphi, l, \theta, o) = \Tv$ in every model
at every point).
\end{theorem}

\begin{proof}
Consider the model $\mathcal{M}_\Uv$ on any frame with
$V(p, l, \theta, o) = \Uv$ for every atom $p$ and every point. By
induction on $\varphi$, using the Kleene tables (each connective returns
$\Uv$ when all inputs are $\Uv$), $\Vh(\varphi, l, \theta, o) = \Uv$ for
every $\varphi$ and every $\frm$. Hence no formula evaluates to $\Tv$
universally. This re-derives, in the indexed setting, the well-known
absence of $K_3$-tautologies under designated value $\Tv$~\cite{kleene1952}.
\end{proof}

\begin{corollary}[Failure of LEM and of contradiction]\label{cor:lem}
$\not\entails p \lor \lnot p$ and $\not\entails \lnot(p \land \lnot p)$.
\end{corollary}

\begin{proof}
Take $\mathcal{M}_\Uv$ as in Theorem~\ref{thm:notaut}: both formulas
evaluate to $\Uv$, not $\Tv$.
\end{proof}

\begin{remark}
Theorem~\ref{thm:notaut} is the price of taking $\Tv$ as the sole
designated value in $K_3$. It does not mean the logic is useless---it
means valid \emph{inferences} (consequences from non-empty $\Gamma$) are
the appropriate notion of validity. For example, $\{p, p \to q\} \entails
q$ holds: by inspection of the implication and conjunction tables, whenever
$\Vh(p) = \Vh(p \to q) = \Tv$, we have $\Vh(q) = \Tv$.
\end{remark}

\section{Worked Example}
\label{sec:example}

Let $L = \{l_{\mathrm{EU}}, l_{\mathrm{US}}\}$, $\Theta = \{\theta_0,
\theta_1\}$, $O = \{o_{\mathrm{Alice}}, o_{\mathrm{Bob}}\}$. Let $p$ be the
atomic proposition ``the package has been delivered'' and $q$ ``the package
has cleared customs.'' A possible base valuation:
\begin{center}
\small
\begin{tabular}{lll|cc}
\toprule
$l$ & $\theta$ & $o$ & $V(p)$ & $V(q)$ \\
\midrule
$l_{\mathrm{EU}}$ & $\theta_0$ & $o_{\mathrm{Alice}}$ & $\Uv$ & $\Fv$ \\
$l_{\mathrm{EU}}$ & $\theta_1$ & $o_{\mathrm{Alice}}$ & $\Tv$ & $\Tv$ \\
$l_{\mathrm{EU}}$ & $\theta_1$ & $o_{\mathrm{Bob}}$   & $\Uv$ & $\Tv$ \\
$l_{\mathrm{US}}$ & $\theta_1$ & $o_{\mathrm{Alice}}$ & $\Fv$ & $\Uv$ \\
\bottomrule
\end{tabular}
\end{center}
Then at $(l_{\mathrm{EU}}, \theta_1, o_{\mathrm{Alice}})$,
$\Vh(p \land q) = \Tv$ (Alice in the EU at $\theta_1$ knows both); at
$(l_{\mathrm{EU}}, \theta_1, o_{\mathrm{Bob}})$, $\Vh(p \land q) = \Uv$
(Bob does not yet know whether the package has been delivered, even though
he knows it has cleared customs); at $(l_{\mathrm{US}}, \theta_1,
o_{\mathrm{Alice}})$, $\Vh(p \land q) = \Fv$ (the proposition is false in
the US frame because $\Vh(p) = \Fv$). Different points yield different
truth values for \emph{the same formula}---this is the design.

\section{Related Work}
\label{sec:related}

The components of LTO-K3 are each well-studied.

\paragraph{Three-valued logics.} $K_3$~\cite{kleene1952} is one of several
three-valued logics with different motivations and different
connectives. {\L}ukasiewicz $\text{\L}_3$~\cite{lukasiewicz1920} differs
from $K_3$ in implication: $\Uv \to_{\text{\L}_3} \Uv = \Tv$, which
yields the tautology $p \to p$ and other classical theorems. Bochvar's
internal $B_3$~\cite{bochvar1981} is ``infectious'': any operation with a
$\Uv$ input returns $\Uv$, even $\Tv \lor \Uv$. Priest's
LP~\cite{priest1979} uses the Kleene tables but designates
$\{\Tv, \Uv\}$, yielding a paraconsistent logic in which contradictions can
be true and LEM holds. Belnap's $\mathbf{FOUR}$~\cite{belnap1977} adds a
fourth value (``both true and false'') and is appropriate when contradictory
information is to be tolerated rather than collapsed to $\Uv$.

We choose $K_3$ with designated value $\Tv$ because the reading of $\Uv$ we
want is purely epistemic (``not yet known'') and information-monotonic
(Proposition~\ref{prop:mono}): resolving an $\Uv$ to a classical value
never overturns a previously determined truth. $\text{\L}_3$ breaks this
property (its implication is not monotone in the obvious sense); LP changes
the designated set; $B_3$ propagates $\Uv$ too aggressively for our use
case; $\mathbf{FOUR}$ conflates uncertainty with inconsistency.

\paragraph{Temporal and modal logics.} Prior~\cite{prior1957} introduced
tense operators; Pnueli~\cite{pnueli1977} introduced LTL for program
verification; Clarke and Emerson~\cite{clarke1981} introduced CTL. These
logics index truth by time (or by states reachable in time) but use
classical two-valued semantics and quantify over states with modal
operators. Kripke~\cite{kripke1963} gave possible-worlds semantics for
modal logic; the entire modal-logic
tradition~\cite{blackburn2001} indexes truth by worlds and uses
$\Box, \Diamond$ to quantify over them. Hybrid
logic~\cite{areces2007hybrid} adds nominals $i$ and the operator $@_i$ to
refer to specific worlds in the object language. LTO-K3 differs by
\emph{not} quantifying over points in the object language at all---there
are no $\Box, \Diamond, \mathsf{F}, \mathsf{G}, K_o, @_i$. Indices live
only in the meta-language.

\paragraph{Epistemic logics.} Hintikka~\cite{hintikka1962} introduced
epistemic modal logic with $K_a$ operators. The textbook treatment is
Fagin et al.~\cite{fagin1995}. Multi-agent epistemic temporal logic
combines these dimensions, but again with two-valued semantics and modal
operators. Three-valued epistemic logics exist but are typically formulated
with modal operators rather than as a flat indexed valuation.

\paragraph{Three-valued spatio-temporal and runtime-verification logics.}
Bauer et al.~\cite{bauer2011rv} developed LTL$_3$, a three-valued LTL
variant for runtime monitoring with truth values true, false, and
inconclusive. Spatial logics~\cite{aiello2007handbook} introduce explicit
spatial indices for reasoning about topological structure. These works are
the closest neighbours of LTO-K3.

\paragraph{Honest positioning.} Relative to the above, LTO-K3 is best read
as a deliberately impoverished system. It has fewer expressive resources
than hybrid logic, fewer than multi-modal epistemic temporal logic, and
fewer than LTL$_3$. What it offers in return is a uniform treatment of
three indices (location, time, observer) as concrete and equally
first-class, and a self-contained semantics that fits on a page. We make
no claim that LTO-K3 supersedes any of these systems; we claim only that it
is a useful \emph{starting point} when their full machinery is not
required.

\section{Limitations}
\label{sec:limitations}

\begin{enumerate}
  \item \textbf{No cross-frame operators.} LTO-K3 cannot express ``$\varphi$
  at $l_1$ implies $\psi$ at $l_2$'' or ``$\varphi$ now will be $\psi$
  later.'' Such statements require modal operators, hybrid nominals, or
  explicit bridging rules, none of which are part of the core.
  \item \textbf{No quantification.} LTO-K3 cannot express
  ``$\varphi$ holds for some observer'' or ``$\varphi$ holds at all
  times.'' These require either propositional quantifiers or
  meta-quantification.
  \item \textbf{Structure on $L, \Theta, O$ is invisible.} Time is not
  ordered, locations are not topologized, observers do not have
  information-state relations. Any such structure must be added by an
  extension and cannot be queried within $\mathcal{L}$.
  \item \textbf{No $\Tv$-tautologies} (Theorem~\ref{thm:notaut}). Validity
  must be reasoned about via non-trivial entailments, not via tautologies.
  Users accustomed to classical-logic tautological reasoning will need to
  adjust.
  \item \textbf{$\Uv$ is undifferentiated.} ``Not yet known'', ``in
  principle unknowable'', and ``not yet defined'' are collapsed into a
  single value. Distinguishing them requires a richer truth-value set
  (e.g.,~\cite{belnap1977}) or a probabilistic / possibilistic layer.
  \item \textbf{No interaction between observers.} There is no notion of
  ``$o_1$ knows that $o_2$ knows $\varphi$''; that requires an explicit
  $K_o$ operator and is intentionally outside scope.
  \item \textbf{The cross-frame undefinedness is a restriction, not a
  feature.} We have framed it as a deliberate design choice, and it is, but
  it is also a real expressive limitation. Users who need cross-frame
  reasoning will need to layer on additional machinery.
\end{enumerate}

\section{Conclusion}
We have presented LTO-K3, a minimal three-valued logic in which every
proposition is evaluated under an explicit triple of location, time, and
observer. The system is a careful packaging of standard
ingredients---Kleene's $K_3$, indexed valuations, designated-value
consequence---into a single self-contained core, together with a small
body of meta-theoretic properties (locality, monotonicity,
conservativity, De~Morgan, absence of $\Tv$-tautologies). We have been
explicit about what is new (the packaging) and what is not (everything
else), and have named the system's limitations. Cross-frame composition,
quantification, and modal operators are deliberate omissions to be
addressed by domain-specific extensions.

\begin{thebibliography}{99}

\bibitem{lukasiewicz1920}
J.~{\L}ukasiewicz, ``O logice tr\'ojwarto\'sciowej'' [On three-valued
logic], \emph{Ruch Filozoficzny}, vol.~5, pp.~169--171, 1920.

\bibitem{kleene1952}
S.~C. Kleene, \emph{Introduction to Metamathematics}, North-Holland,
Amsterdam, 1952.

\bibitem{prior1957}
A.~N. Prior, \emph{Time and Modality}, Oxford University Press, Oxford,
1957.

\bibitem{hintikka1962}
J.~Hintikka, \emph{Knowledge and Belief: An Introduction to the Logic of
the Two Notions}, Cornell University Press, Ithaca, NY, 1962.

\bibitem{kripke1963}
S.~A. Kripke, ``Semantical considerations on modal logic,''
\emph{Acta Philosophica Fennica}, vol.~16, pp.~83--94, 1963.

\bibitem{belnap1977}
N.~D. Belnap, ``A useful four-valued logic,'' in J.~M. Dunn and
G.~Epstein, eds., \emph{Modern Uses of Multiple-Valued Logic}, Reidel,
Dordrecht, pp.~5--37, 1977.

\bibitem{pnueli1977}
A.~Pnueli, ``The temporal logic of programs,'' in \emph{Proc.\ 18th
Annual Symposium on Foundations of Computer Science (FOCS)}, IEEE,
pp.~46--57, 1977.

\bibitem{priest1979}
G.~Priest, ``The logic of paradox,'' \emph{Journal of Philosophical
Logic}, vol.~8, no.~1, pp.~219--241, 1979.

\bibitem{bochvar1981}
D.~A. Bochvar, ``On a three-valued logical calculus and its application to
the analysis of the paradoxes of the classical extended functional
calculus,'' translated by M.~Bergmann, \emph{History and Philosophy of
Logic}, vol.~2, pp.~87--112, 1981 (original Russian 1938).

\bibitem{clarke1981}
E.~M. Clarke and E.~A. Emerson, ``Design and synthesis of synchronization
skeletons using branching-time temporal logic,'' in \emph{Logics of
Programs}, LNCS 131, Springer, pp.~52--71, 1981.

\bibitem{fagin1995}
R.~Fagin, J.~Y. Halpern, Y.~Moses, and M.~Y. Vardi, \emph{Reasoning About
Knowledge}, MIT Press, Cambridge, MA, 1995.

\bibitem{blackburn2001}
P.~Blackburn, M.~de~Rijke, and Y.~Venema, \emph{Modal Logic}, Cambridge
University Press, Cambridge, 2001.

\bibitem{aiello2007handbook}
M.~Aiello, I.~Pratt-Hartmann, and J.~van~Benthem, eds., \emph{Handbook of
Spatial Logics}, Springer, Dordrecht, 2007.

\bibitem{areces2007hybrid}
C.~Areces and B.~ten~Cate, ``Hybrid logics,'' in P.~Blackburn,
J.~van~Benthem, and F.~Wolter, eds., \emph{Handbook of Modal Logic},
Elsevier, pp.~821--868, 2007.

\bibitem{bauer2011rv}
A.~Bauer, M.~Leucker, and C.~Schallhart, ``Runtime verification for
LTL and TLTL,'' \emph{ACM Transactions on Software Engineering and
Methodology}, vol.~20, no.~4, art.~14, 2011.

\end{thebibliography}

\end{document}
